\chapter{Schlussfolgerungen, Limitationen und Ausblick} % 5-10%
\label{ch:Fazit}

% In der Diskussion/im Fazit bzw. Schlussteil der Arbeit
%• gilt es, die empirischen Ergebnisse mit dem bereits vorgestellten Forschungsstand in Beziehung zu setzen. Wo decken oder widersprechen sich Ergebnisse? Welche Fragen sind offen geblieben?
%• Welche weiterführenden Fragen/Forschungsbedarf oder ggf. Handlungsempfehlungen für die Praxis lassen sich formulieren?

% FF: Wie sinnvoll ist die Einführung eines modernen Dashboards im Hinblick auf das Ziel, den Legacy-Testpoolmonitoring-Prozess effizienter und nutzerfreundlicher zu gestalten?


% Aufgreifen der Zielsetzung
\textcolor{red}{TODO: EInleitungssatz -} In einer Welt voller Altsysteme, hatte ich xy vor ...\\
Um einen Beitrag zu nachhaltigen und effizienten Systemen beizutragen, sollte mit dieser Forschungsarbeit herausgestellt werden, welchen Beitrag moderne Dashboards zur Modernisierung beitragen können eine geeignete Lösung darstellen. Dafür wurde anhand eines Legacy-Testpools ein Altsystem derer Ergebnisanzeige als Grundlage für ein neues Dashboard aufgegriffen. Ziel war es herauszustellen, wie sinnvoll dessen Einführung ist, um den Monitoring-Prozess effizienter und nutzerfreundlicher zu gestalten.

% Zusammenfassung
\section{Zusammenfassung der wichtigsten Ergebnisse}

Hierfür wurde eine Nutzerstudie bestehend aus einem Usability-Test und einer dazugehörigen Nutzerbefragung durchgeführt. Während die Nutzerbefragung auf den Prinzipien der ISO 9241-110 basiert, wurden innerhalb des Usability-Tests die nachträglich abgefragten Eigenschaften eingebaut. Auf dieser Grundlage konnten die beide Systeme anhand der drei Usability-Dimensionen Effektivität, Effizienz und Nutzerzufriedenheit verglichen werden.

Anhand dieser Forschungen konnten einige Möglichkeiten festgestellt werden.

Mithilfe der gewählten Dashboard-Lösung, ist ein zeitlich effizienteres Monitoring möglich als zuvor. Zwar wies der Usability-Test an dem Dashboard eine höhere Fehlerquote auf, als an dem Altsystem, jedoch liegt dies an der exzellenten Kenntnis der Nutzer des Altsystems sowie der Tatsache, dass die neue Lösung erst mit der Teilnahme an der Studie vorgestellt wurde. Allerdings konnte das Dashboard in den relevanten Stellen des Tests durch seine verbesserte Effektivität überzeugen und somit eine knapp 50\% bessere Effizienz vorweisen.

Anhand der regen Teilnahme an der Feedbackrunde war festzustellen, dass durchaus Verbesserungspotenzial an dem neuen Prototypen besteht. Jedoch wird damit auch Interesse auf Nutzerbasis gezeigt, dieses Projekt weiter zu fördern. Insgesamt sind die Nutzer auch viel zufriedener mit der Umsetzung und Nutzbarkeit des Dashboards, als mit der gewachsenen Altlösung.

% Einordnung in den aktuellen Forschungsstand
\section{Einordnung in den aktuellen Forschungsstand}

Die Dashboard-Lösung bildet eine nicht-invasive, risikoarme Modernisierungsmethode. Es hilft Anwendern Monitoring-Prozesse effizienter zu gestalten, da es die alten Strukturen in eine nutzerfreundliche Oberfläche verpackt. Da diese Dashboard-Lösung jedoch keine bestehenden Strukturen entfernt, entsteht auch keine Hürde für ältere Mitarbeiter, welche potenziell Probleme mit der Umstellung haben.

Weiterhin bietet es sich an, eine Ergebnisanzeige so zu modernisieren, um vergessene oder geduldete Altlasten aufzudecken. Da es sich teilweise schwierig gestaltet, diese durch die Modernisierungsmaßnahme automatisiert auszulesen, bietet der Prozess eine geeignete Gelegenheit, die Altlasten neu zu entdecken und aufzuarbeiten.

Es bietet darüber hinaus die Gelegenheit, auf das Thema der Modernisierung aufmerksam zu machen, wie diese beispielsweise eine zur Optimierung von Prozessen eingesetzt werden kann. Das Dashboard kann somit als eine geeignete Modernisierungsmethode für Legacy-Anwendungen mit Ergebnisanzeige eingeordnet werden.

% Grenzen + mögliche Weiterentwicklungen
\section{Grenzen des entwickelten Prototypen und mögliche Weiterentwicklungen}

Da die Dashboard-Lösung bislang ein Projekt prototypischen Charakters ist, kann dieses Projekt noch stark ausgeweitet werden. Neben dem Ausbau auf den restlichen Testpool, kann das bereits genannte Nutzerfeedback zur Optimierung angewendet werden. Dieses liefert hilfreiche Hinweise bezüglich dem bestehenden Workflow sowie Formatierung, Größe und Kapselung zusammengehöriger Informationen, um den Übergang auf das neue System zu erleichtern.

Bisher bietet das Dashboard zudem rein informell betrachtet keinen Mehrwert. Da die Daten unabhängig vom bestehenden System ausgewertet und aufgearbeitet werden, sind vorgeschlagene Features wie die Darstellung der Fehlerwerte der letzten zwei Wochen in einem Verlaufsdiagramm möglich. Da das Dashboard jeden Tag überwacht wird, ist ein direkter Vergleichswert zum Vortag denkbar. Dieser zeigt zusätzlich zur Fehlerzahl, ob und wie viele Fehler mehr oder weniger als am Vortag vorliegen.

Da das System keine Veränderung am Altsystem vornimmt, kann diese Lösung ausschließlich als erster Schritt des Modernisierungsprozesses angesehen werden. Es hilft Altlasten zu erkennen und fachfremden Personen den Einstieg in das System zu erleichtern. Ein nächster Schritt wäre, diese Erkenntnisse zur Umstrukturierung und Optimierung oder zur Umstrukturierung in Kombination mit einem Rehosting auf zum Beispiel ein standardisiertes, modernes Framework zu nutzen.

% Ausblick + Empfehlungen für weitergehende Forschung
\section{Ausblick und Empfehlungen für weitergehende Forschung}

Dieses Prototyp-Projekt kann in naher Zukunft mit den aufgezählten Features und vorgeschlagenen Verbesserungen ausgebaut und optimiert werden.
Dabei können die Erkenntnisse aus dieser Forschungsarbeit eine Grundlage für beispielsweise folgende Aufbaustudien bieten:

Aufgrund der begrenzten Kapazitäten ist eine Wiederholung des Versuchs mit einer höheren Teilnehmerzahl, oder  einem größeren Umfang in Funktionalität und Aufgaben denkbar, um die Ergebnisse mit einer größeren Teststärke validieren zu können.
Mit mehr Teilnehmern einer Studie kann auf eine breitere Fächerung sowie gleichmäßigere Verteilung der Nutzer in Erfahrungslevel, Expertise sowie Altersklasse geachtet werden. Dies ermöglicht, Rückschlüsse auf die verschiedenen Variablen ziehen zu können. Dieser Ansatz kann Erkenntnisse darauf bieten, wie Modernisierungsmaßnahmen aufgefasst werden oder wie diese in den verschiedenen Nutzergruppen angegangen werden sollten.

Eine weitere Möglichkeit stellt die Untersuchung tiefergreifender Modernisierungsmaßnahmen wie der Umstrukturierung oder des Rehostings dar. Es könnten verschiedene Ansätze getestet oder unterschiedliche standardisierte Frameworks miteinander verglichen werden, um eine gute Lösung für das jeweilige System finden zu können.

\textcolor{red}{TODO: brauche ich einen weiteren Abschlusssatz?}

% -> was ist meine Empfehlung (potenzielle neue Arbeiten) die ich gebe, damit diese Forschung in der Zukunft eingesetzt werden kann
